%%
%% This is file `sample-MapleTransactions.tex',
%% modified from sample-acmsmall.tex
%% which is part of the acmart.cls
%% template collection.
%%
%% The original source files were:
%%
%% samples.dtx  (with options: `acmsmall')
%%
%% IMPORTANT NOTICE:
%%
%% For the copyright see the source file.
%%
%%
%% For distribution of the original source see the terms
%% for copying and modification in the file samples.dtx.
%%
%% This generated file may be distributed as long as the
%% original source files, as listed above, are part of the
%% same distribution. (The sources need not necessarily be
%% in the same archive or directory.)
%%
%% The first command in your LaTeX source must be the \documentclass command.
\documentclass[anonymous,review,acmsmall,screen]{MapleTrans}
%% NOTE that a single column version is used in this journal. %%
%% useful options are
%% \documentclass[anonymous,acmsmall,screen]
%% which anonymizes the article.
%%
%% The "screen" option puts colors
%% on hyperlinks for ease of
%% distinguishing onscreen.
%%
%% On packages:
%% To ensure 100% compatibility, please check the white list of
%% approved LaTeX packages to be used with the Master Article Template at
%% https://www.acm.org/publications/taps/whitelist-of-latex-packages
%% before creating your document. The white list page provides
%% information on how to submit additional LaTeX packages for
%% review and adoption.

%% The following are useful packages
%\usepackage[backref=true, style=numeric]{biblatex}
%\addbibresource{sample-base.bib}
\usepackage[normalem]{ulem}
\newcommand{\strikethrough}[1]{\sout{#1}}
\usepackage{enumitem}
%\usepackage{hyperref} % already included?
%\usepackage{mathtools} % already included in MapleTrans.cls
\usepackage{etoolbox}
%\usepackage{euler} % incompatible
%
% The package breqn.sty is
% incompatible with acmart.cls
% and therefore also with MapleTrans.cls
% unless one loads "lineno.sty"
% and then redefines the "dmath"
% environment.
%
\usepackage{breqn}
\usepackage{lineno}
% Overleaf doesn't like the \begin stuff below, but it's fine
%\BeforeBeginEnvironment{dmath}{\begin{nolinenumbers}}%
%\AfterEndEnvironment{dmath}{\end{nolinenumbers}}
\usepackage{maple}
% Other useful packages
\usepackage{xfrac}
\usepackage{xcolor}
\definecolor{mygreen}{rgb}{0,0.6,0}
\definecolor{mygray}{rgb}{0.5,0.5,0.5}
\definecolor{mymauve}{rgb}{0.58,0,0.82}
\definecolor{altblue}{rgb}{0.0,0.6,1.0}
\definecolor{lstbg}{cmyk}{0.05, 0.01, 0, 0}
\definecolor{morebluish}{cmyk}{0.06,0.04,0,0}
% Enclose Maple code (short
% passages only) in the lstlisting
% environment with language set
% to "maple", as follows.
% (uses colours defined above)
\usepackage{listings}
\input{listings-maple-definition.sty}
\lstset{
backgroundcolor=\color{lstbg},
  % choose the background color; you must add \usepackage{color} or \usepackage{xcolor}
basicstyle=\small\ttfamily,language=maple
}

%% Rights management information.  This information is sent to you
%% when you complete the rights form.  These commands have SAMPLE
%% values in them; it is your responsibility as an author to replace
%% the commands and values with those provided to you when you
%% complete the rights form.
\setcopyright{acmlicensed}
\copyrightyear{2021}
\acmYear{2021}
% For editorial use only:
% parameter V = Volume number
% parameter I = Issue number
% parameter XXXXX = article number
\acmDOI{10.5206/mt.vViI.XXXXX}



%%
%% These commands are for a JOURNAL article.
\acmJournal{MAPLETRANS}
\acmVolume{X}
\acmNumber{Y}
\acmArticle{Z}
\acmMonth{0}

% User macros
\newcommand{\BibTeX}{\texttt{BibTeX}}


%%
%% end of the preamble, start of the body of the document source.
\begin{document}

%%
%% The "title" command has an optional parameter,
%% allowing the author to define a "short title" to be used in page headers.
\title[MT Template 1.0]{Maple Transactions \LaTeX{} Template Version 1.0}

%%
%% The "author" command and its associated commands are used to define
%% the authors and their affiliations.
%% Of note is the shared affiliation of the first two authors, and the
%% "authornote" and "authornotemark" commands
%% used to denote shared contribution to the research.
\author{Ben Trovato}
\authornote{Both authors contributed equally to this research.}
\email{trovato@corporation.com}
\orcid{1234-5678-9012}
\author{G.K.M. Tobin}
\authornotemark[1]
\email{webmaster@marysville-ohio.com}
\affiliation{%
  \institution{Institute for Clarity in Documentation}
  \streetaddress{P.O. Box 1212}
  \city{Dublin}
  \state{Ohio}
  \country{USA}
  \postcode{43017-6221}
}

\author{Lars Th{\o}rv{\"a}ld}
\affiliation{%
  \institution{The Th{\o}rv{\"a}ld Group}
  \streetaddress{1 Th{\o}rv{\"a}ld Circle}
  \city{Hekla}
  \country{Iceland}}
\email{larst@affiliation.org}

\author{Valerie B\'eranger}
\affiliation{%
  \institution{Inria Paris-Rocquencourt}
  \city{Rocquencourt}
  \country{France}
}

\author{Aparna Patel}
\affiliation{%
 \institution{Rajiv Gandhi University}
 \streetaddress{Rono-Hills}
 \city{Doimukh}
 \state{Arunachal Pradesh}
 \country{India}}

\author{Huifen Chan}
\affiliation{%
  \institution{Tsinghua University}
  \streetaddress{30 Shuangqing Rd}
  \city{Haidian Qu}
  \state{Beijing Shi}
  \country{China}}

\author{Charles Palmer}
\affiliation{%
  \institution{Palmer Research Laboratories}
  \streetaddress{8600 Datapoint Drive}
  \city{San Antonio}
  \state{Texas}
  \country{USA}
  \postcode{78229}}
\email{cpalmer@prl.com}

\author{John Smith}
\affiliation{%
  \institution{The Th{\o}rv{\"a}ld Group}
  \streetaddress{1 Th{\o}rv{\"a}ld Circle}
  \city{Hekla}
  \country{Iceland}}
\email{jsmith@affiliation.org}

\author{Julius P. Kumquat}
\affiliation{%
  \institution{The Kumquat Consortium}
  \city{New York}
  \country{USA}}
\email{jpkumquat@consortium.net}

%%
%% By default, the full list of authors will be used in the page
%% headers. Often, this list is too long, and will overlap
%% other information printed in the page headers. This command allows
%% the author to define a more concise list
%% of authors' names for this purpose.
\renewcommand{\shortauthors}{Trovato and Tobin, et al.}

%%
%% The abstract is a short summary of the work to be presented in the
%% article.
\begin{abstract}
\textcolor{red}{ \textbf{Abstract}.
This template is to show how to use MapleTrans.cls, which is a \LaTeX\ class file derived from acmart.cls. Differences from the acmart.cls will be marked; here we use the color red to do so.  One difference is that Maple Transactions would like you to include the word ``Abstract" in bold at the beginning of your abstract.}
 A clear and well-documented \LaTeX\ document is presented as an
  article formatted for publication by ACM in a conference proceedings
  or journal publication. Based on the ``acmart'' document class, this
  article presents and explains many of the common variations, as well
  as many of the formatting elements an author may use in the
  preparation of the documentation of their work.
\end{abstract}

%%
%% The code below is generated by the tool at http://dl.acm.org/ccs.cfm.
%% Please copy and paste the code instead of the example below.
%%
\begin{CCSXML}
<ccs2012>
 <concept>
  <concept_id>10010520.10010553.10010562</concept_id>
  <concept_desc>Computer systems organization~Embedded systems</concept_desc>
  <concept_significance>500</concept_significance>
 </concept>
 <concept>
  <concept_id>10010520.10010575.10010755</concept_id>
  <concept_desc>Computer systems organization~Redundancy</concept_desc>
  <concept_significance>300</concept_significance>
 </concept>
 <concept>
  <concept_id>10010520.10010553.10010554</concept_id>
  <concept_desc>Computer systems organization~Robotics</concept_desc>
  <concept_significance>100</concept_significance>
 </concept>
 <concept>
  <concept_id>10003033.10003083.10003095</concept_id>
  <concept_desc>Networks~Network reliability</concept_desc>
  <concept_significance>100</concept_significance>
 </concept>
</ccs2012>
\end{CCSXML}

\ccsdesc[500]{Computer systems organization~Embedded systems}
\ccsdesc[300]{Computer systems organization~Redundancy}
\ccsdesc{Computer systems organization~Robotics}
\ccsdesc[100]{Networks~Network reliability}

%%
%% Keywords. The author(s) should pick words that accurately describe
%% the work being presented. Separate the keywords with commas.
\keywords{datasets, neural networks, gaze detection, text tagging}


%%
%% This command processes the author and affiliation and title
%% information and builds the first part of the formatted document.
\maketitle

\section{Introduction}
\textcolor{red}{The editors of Maple Transactions chose to base the style file for the journal on acmart.cls, for several reasons, including the following:
\begin{enumerate}
    \item Many authors for Maple Transactions will already be familiar with acmart.cls, so learning to use this one should be straightforward.
    \item The acmart.cls is well-designed and has several desirable features, including an easy switch ``anonymous" for anonymous author/anonymous reviewer refereeing (also called ``double-blind'' refereeing) which Maple Transactions has adopted.
    \item Permission to modify the acmart.cls is granted as part of the LaTex Public Project license \href{https://www.latex-project.org/lppl/}{https://www.latex-project.org/lppl/}
    \item An important point about those permissions: please do not send bug reports about this template to ACM!  Please send them instead to \textcolor{blue}{\texttt{rcorless@uwo.ca}}.
\end{enumerate}
}
ACM's consolidated article template, introduced in 2017, provides a
consistent \LaTeX\ style for use across ACM publications, and
incorporates accessibility and metadata-extraction functionality
necessary for future Digital Library endeavors. Numerous ACM and
SIG-specific \LaTeX\ templates have been examined, and their unique
features incorporated into this single new template.

If you are new to publishing with ACM, this document is a valuable
guide to the process of preparing your work for publication. If you
have published with ACM before, this document provides insight and
instruction into more recent changes to the article template.

The ``\verb|acmart|'' document class can be used to prepare articles
for any ACM publication --- conference or journal, and for any stage
of publication, from review to final ``camera-ready'' copy, to the
author's own version, with {\itshape very} few changes to the source.



\section{Template Overview}
As noted in the introduction, the ``\verb|acmart|'' document class can
be used to prepare many different kinds of documentation --- a
double-blind initial submission of a full-length technical paper, a
two-page SIGGRAPH Emerging Technologies abstract, a ``camera-ready''
journal article, a SIGCHI Extended Abstract, and more --- all by
selecting the appropriate {\itshape template style} and {\itshape
  template parameters}.

This document will explain the major features of the document
class. For further information, the {\itshape \LaTeX\ User's Guide} is
available from
\url{https://www.acm.org/publications/proceedings-template}.

\subsection{Template Styles}

The primary parameter given to the ``\verb|acmart|'' document class is
the {\itshape template style} which corresponds to the kind of publication
or SIG publishing the work. This parameter is enclosed in square
brackets and is a part of the {\verb|documentclass|} command:
\begin{verbatim}
  \documentclass[STYLE]{acmart}
\end{verbatim}
\textcolor{red}{To use MapleTrans, use either of the following:}
\begin{verbatim}
    \documentclass[acmsmall,screen]{MapleTrans}
\end{verbatim}
\textcolor{red}{or}
\begin{verbatim}
    \documentclass[anonymous,review,acmsmall,screen]{MapleTrans}
\end{verbatim}
\textcolor{red}{The anonymous option is helpful for anonymous author/anonymous referee reviewing, which Maple Transactions has adopted.}
\textcolor{red}{The ``review'' option puts line numbers everywhere.  These may not \textsl{quite} line up with every line, but they are quite useful for communicating feedback to authors.}


Journals use one of three template styles. All but three ACM journals
use the {\verb|acmsmall|} template style:
\begin{itemize}
\item {\verb|acmsmall|}: The default journal template style.
\item {\verb|acmlarge|}: Used by JOCCH and TAP.
\item {\verb|acmtog|}: Used by TOG.
\end{itemize}

The majority of conference proceedings documentation will use the {\verb|acmconf|} template style.
\begin{itemize}
\item {\verb|acmconf|}: The default proceedings template style.
\item{\verb|sigchi|}: Used for SIGCHI conference articles.
\item{\verb|sigchi-a|}: Used for SIGCHI ``Extended Abstract'' articles.
\item{\verb|sigplan|}: Used for SIGPLAN conference articles.
\end{itemize}

\subsection{Template Parameters}

In addition to specifying the {\itshape template style} to be used in
formatting your work, there are a number of {\itshape template parameters}
which modify some part of the applied template style. A complete list
of these parameters can be found in the {\itshape \LaTeX\ User's Guide.}

Frequently-used parameters, or combinations of parameters, include:
\begin{itemize}
\item {\verb|anonymous,review|}: Suitable for a ``double-blind''
  conference submission. Anonymizes the work and includes line
  numbers. Use with the \verb|\acmSubmissionID| command to print the
  submission's unique ID on each page of the work.
\item{\verb|authorversion|}: Produces a version of the work suitable
  for posting by the author.
\item{\verb|screen|}: Produces colored hyperlinks.
\end{itemize}

This document uses the following string as the first command in the
source file:
\begin{verbatim}
\documentclass[acmsmall]{acmart}
\end{verbatim}

\section{An important incompatibility issue}
\textcolor{red}{The acmart.cls file on which the Maple Transactions style MapleTrans.cls is based is incompatible with the breqn.sty file, which is a delicate style file with a powerful tool that is desirable for use with computer algebra, namely automatic line breaking for multiline equations.
Currently, exporting a Maple document to LaTeX requires breqn.sty and maple.sty which depends on breqn.sty.
At this time, this incompatibility has not been fully resolved.  Instead of using the \texttt{maple.sty} supplied with Maple 2021.1 \emph{or earlier}, one needs to use the \texttt{maple.sty} file attached to this template.  This will become the \texttt{maple.sty} file distributed with Maple 2021.2.}

\textcolor{red}{In order that breqn's dmath work with acmart.cls or with MapleTrans.cls, one must be careful and load the packages in the order specified below, and use the specified command from the ``lineno" package to adjust the dmath environment so that it won't cause the pdf generator of Overleaf to hang.  We hope this situation is temporary. Indeed, the issue (in Overleaf) seems to have been fixed now, and in this template those lines have now been commented out.
[Remark: the separation of the \texttt{\textbackslash begin}
and the
\texttt{\textbackslash end} as shown below (commented out) caused some \LaTeX\ front-ends to complain; for instance, Overleaf's syntax-checker does not like that separation.  It is, however, legal \LaTeX\ and does what is intended, namely prevent the pdf generator from hanging, so you may ignore the warning that happens if you uncomment the lines.]
}
\begin{verbatim}
\usepackage{etoolbox}
\usepackage{breqn}
\usepackage{lineno}
%\BeforeBeginEnvironment{dmath}{\begin{nolinenumbers}}%
%\AfterEndEnvironment{dmath}{\end{nolinenumbers}}
\usepackage{maple}
\end{verbatim}



\section{Modifications}

\textcolor{red}{The acmart.cls is quite restrictive with respect to modifications: }
Modifying the template --- including but not limited to: adjusting
margins, typeface sizes, line spacing, paragraph and list definitions,
and the use of the \verb|\vspace| command to manually adjust the
vertical spacing between elements of your work --- \strikethrough{is not allowed}.

\strikethrough{\bfseries Your document will be returned to you for revision if
  modifications are discovered.}
\textcolor{red}{Maple Transactions, however, is much less restrictive; many modifications are permitted.  The editors would like to keep a more-or-less uniform look to the online journal, but if the author feels strongly about something, the editors are likely to be sympathetic.}


\section{Typefaces}

The ``\verb|acmart|'' document class requires the use of the
``Libertine'' typeface family. Your \TeX\ installation should include
this set of packages. Please do not substitute other typefaces. The
``\verb|lmodern|'' and ``\verb|ltimes|'' packages should not be used,
as they will override the built-in typeface families.
\textcolor{red}{Again, Maple Transactions is less restrictive, and alternative typefaces (within reason) will be permitted.}


\section{Title Information}

The title of your work should use capital letters appropriately -
\url{https://capitalizemytitle.com/} has useful rules for
capitalization. Use the {\verb|title|} command to define the title of
your work. If your work has a subtitle, define it with the
{\verb|subtitle|} command.  Do not insert line breaks in your title.

If your title is lengthy, you must define a short version to be used
in the page headers, to prevent overlapping text. The \verb|title|
command has a ``short title'' parameter:
\begin{verbatim}
  \title[short title]{full title}
\end{verbatim}

\section{Authors and Affiliations}

Each author must be defined separately for accurate metadata
identification. Multiple authors may share one affiliation. Authors'
names should not be abbreviated; use full first names wherever
possible. Include authors' e-mail addresses whenever possible.

Grouping authors' names or e-mail addresses, or providing an ``e-mail
alias,'' as shown below, is not acceptable:
\begin{verbatim}
  \author{Brooke Aster, David Mehldau}
  \email{dave,judy,steve@university.edu}
  \email{firstname.lastname@phillips.org}
\end{verbatim}

The \verb|authornote| and \verb|authornotemark| commands allow a note
to apply to multiple authors --- for example, if the first two authors
of an article contributed equally to the work.

If your author list is lengthy, you must define a shortened version of
the list of authors to be used in the page headers, to prevent
overlapping text. The following command should be placed just after
the last \verb|\author{}| definition:
\begin{verbatim}
  \renewcommand{\shortauthors}{McCartney, et al.}
\end{verbatim}
Omitting this command will force the use of a concatenated list of all
of the authors' names, which may result in overlapping text in the
page headers.

The article template's documentation, available at
\url{https://www.acm.org/publications/proceedings-template}, has a
complete explanation of these commands and tips for their effective
use.

Note that authors' addresses are mandatory for journal articles.

\section{Rights Information}

Authors of any work published by \strikethrough{ACM}\textcolor{red}{Maple Transactions} will need to complete a rights
form. Depending on the kind of work, and the rights management choice
made by the author, this may be copyright transfer, permission,
license, or an OA (open access) agreement.

Regardless of the rights management choice, the author will receive a
copy of the completed rights form once it has been submitted. This
form contains \LaTeX\ commands that must be copied into the source
document. When the document source is compiled, these commands and
their parameters add formatted text to several areas of the final
document:
\begin{itemize}
\item the ``ACM Reference Format'' text on the first page.
\item the ``rights management'' text on the first page.
\item \strikethrough{the conference information in the page header(s).}
\end{itemize}

Rights information is unique to the work; if you are preparing several
works for an event, make sure to use the correct set of commands with
each of the works.

\strikethrough{The ACM Reference Format text is required for all articles over one
page in length, and is optional for one-page articles (abstracts)}.

\section{CCS Concepts and User-Defined Keywords}

Two elements of the \strikethrough{``acmart''}\textcolor{red}{MapleTrans} document class provide powerful
taxonomic tools for you to help readers find your work in an online
search.

The ACM Computing Classification System ---
\url{https://www.acm.org/publications/class-2012} --- is a set of
classifiers and concepts that describe the computing
discipline. Authors can select entries from this classification
system, via \url{https://dl.acm.org/ccs/ccs.cfm}, and generate the
commands to be included in the \LaTeX\ source.

User-defined keywords are a comma-separated list of words and phrases
of the authors' choosing, providing a more flexible way of describing
the research being presented.

CCS concepts and user-defined keywords are required for for all
articles over two pages in length, and are optional for one- and
two-page articles (or abstracts).

\section{Sectioning Commands}

Your work should use standard \LaTeX\ sectioning commands:
\verb|section|, \verb|subsection|, \verb|subsubsection|, and
\verb|paragraph|. They should be numbered; do not remove the numbering
from the commands.

Simulating a sectioning command by setting the first word or words of
a paragraph in boldface or italicized text is {\bfseries not allowed.}

\medskip\par
\textcolor{red}{Because the hyperref package is included, several very useful commands are immediately available.}
\begin{enumerate}
    \item \textcolor{red}{href and url for including links to external websites}
    \item \textcolor{red}{hypertarget---a utility for linking within the document.}
\end{enumerate}

\section{Tables}

The ``\verb|acmart|'' document class \textcolor{red}{and the MapleTrans class} include\strikethrough{s} the ``\verb|booktabs|''
package --- \url{https://ctan.org/pkg/booktabs} --- for preparing
high-quality tables.

Table captions are placed {\itshape above} the table.

Because tables cannot be split across pages, the best placement for
them is typically the top of the page nearest their initial cite.  To
ensure this proper ``floating'' placement of tables, use the
environment \textbf{table} to enclose the table's contents and the
table caption.  The contents of the table itself must go in the
\textbf{tabular} environment, to be aligned properly in rows and
columns, with the desired horizontal and vertical rules.  Again,
detailed instructions on \textbf{tabular} material are found in the
\textit{\LaTeX\ User's Guide}.

Immediately following this sentence is the point at which
Table~\ref{tab:freq} is included in the input file; compare the
placement of the table here with the table in the printed output of
this document.

\begin{table}
  \caption{Frequency of Special Characters}
  \label{tab:freq}
  \begin{tabular}{ccl}
    \toprule
    Non-English or Math&Frequency&Comments\\
    \midrule
    \O & 1 in 1,000& For Swedish names\\
    $\pi$ & 1 in 5& Common in math\\
    \$ & 4 in 5 & Used in business\\
    $\Psi^2_1$ & 1 in 40,000& Unexplained usage\\
  \bottomrule
\end{tabular}
\end{table}

To set a wider table, which takes up the whole width of the page's
live area, use the environment \textbf{table*} to enclose the table's
contents and the table caption.  As with a single-column table, this
wide table will ``float'' to a location deemed more
desirable. Immediately following this sentence is the point at which
Table~\ref{tab:commands} is included in the input file; again, it is
instructive to compare the placement of the table here with the table
in the printed output of this document.

\begin{table*}
  \caption{Some Typical Commands}
  \label{tab:commands}
  \begin{tabular}{ccl}
    \toprule
    Command &A Number & Comments\\
    \midrule
    \texttt{{\char'134}author} & 100& Author \\
    \texttt{{\char'134}table}& 300 & For tables\\
    \texttt{{\char'134}table*}& 400& For wider tables\\
    \bottomrule
  \end{tabular}
\end{table*}

Always use midrule to separate table header rows from data rows, and
use it only for this purpose. This enables assistive technologies to
recognise table headers and support their users in navigating tables
more easily.

\section{Math Equations}
You may want to display math equations in three distinct styles:
inline, numbered or non-numbered display.  Each of the three are
discussed in the next sections.

\subsection{Inline (In-text) Equations}
A formula that appears in the running text is called an inline or
in-text formula.  It is produced by the \textbf{math} environment,
which can be invoked with the usual
\texttt{{\char'134}begin\,\ldots{\char'134}end} construction or with
the short form \texttt{\$\,\ldots\$}. You can use any of the symbols
and structures, from $\alpha$ to $\omega$, available in
\LaTeX~\cite{Lamport:LaTeX}; this section will simply show a few
examples of in-text equations in context. Notice how this equation:
\begin{math}
  \lim_{n\rightarrow \infty}x=0
\end{math},
set here in in-line math style, looks slightly different when
set in display style.  (See next section).

\subsection{Display Equations}
A numbered display equation---one set off by vertical space from the
text and centered horizontally---is produced by the \textbf{equation}
environment. An unnumbered display equation is produced by the
\textbf{displaymath} environment.

Again, in either environment, you can use any of the symbols and
structures available in \LaTeX\@; this section will just give a couple
of examples of display equations in context.  First, consider the
equation, shown as an inline equation above:
\begin{equation}
  \lim_{n\rightarrow \infty}x=0
\end{equation}
Notice how it is formatted somewhat differently in
the \textbf{displaymath}
environment.  Now, we'll enter an unnumbered equation:
\begin{displaymath}
  \sum_{i=0}^{\infty} x + 1
\end{displaymath}
and follow it with another numbered equation:
\begin{equation}
  \sum_{i=0}^{\infty}x_i=\int_{0}^{\pi+2} f
\end{equation}
just to demonstrate \LaTeX's able handling of numbering.
\section{Including Maple or other Code}
\textcolor{red}{
Long Maple programs should not be included in the text; provide them instead in the supplementary material (ideally, a link to a working Maple workbook with all code and commands).  Short bits of code, intended to be read as part of the text, are very welcome.
}

\textcolor{red}{This can be done in two ways: first by use of the \texttt{lstlisting} environment.  The following listing (of a meaningless code fragment) is formatted and displayed in Listing~\ref{list:ajrec}.
}
\textcolor{red}{In the preamble, put}
\begin{verbatim}
\usepackage{xcolor}
\definecolor{mygreen}{rgb}{0,0.6,0}
\definecolor{mygray}{rgb}{0.5,0.5,0.5}
\definecolor{mymauve}{rgb}{0.58,0,0.82}
\definecolor{altblue}{rgb}{0.0,0.6,1.0}
\definecolor{lstbg}{cmyk}{0.05, 0.01, 0, 0}
\definecolor{morebluish}{cmyk}{0.06,0.04,0,0}
% Enclose Maple code (short
% passages only) in the lstlisting
% environment with language set
% to "maple", as follows.
% (uses colours defined above)
\usepackage{listings}
\input{listings-maple-definition.sty}
\lstset{
backgroundcolor=\color{lstbg},
  % choose the background color; you must add \usepackage{color} or \usepackage{xcolor}
basicstyle=\small\ttfamily,language=maple
}
\end{verbatim}
\textcolor{red}{and then in the text you may include your code by using commands like the following.}
\begin{verbatim}
\begin{lstlisting}[caption={A reasonably efficient implementation
                            of a recurrence relation}
                \label{list:ajrec}]
  a := proc(n, c);
         if not n::posint then
           return 'procname'(args)
         end if;
         return rememberedA(n,c);
       end proc:
\end{lstlisting}
\end{verbatim}

\begin{lstlisting}[caption={A reasonably efficient implementation of a recurrence relation} \label{list:ajrec}]
a := proc(n, c);
  if not n::posint then
    return 'procname'(args)
  end if;
  return rememberedA(n,c);
end proc:
\end{lstlisting}
\textcolor{red}{
The \texttt{lstlisting} environment can be used similarly for other languages, such as Python or Matlab.  See the documentation of the listings package for examples of use for other languages.
}
\medskip\par
\textcolor{red}{
The second method to include (short) pieces of Maple in the paper is to use output of Maple's
``Export to LaTeX'' command in the file.  This uses the \texttt{maple.sty} file discussed earlier.  With that style file, the following input produces the output below it.
}
\begin{verbatim}
\mapleinput
{$ \displaystyle \texttt{>\,} \mathit{MGF} \coloneqq (y ,c)\,\rightarrow
\,\mathit{evalf} (y -\mathit{Sum} (\frac{a (j ,c)}{y^{2\cdot j -1}},
j =1..\mathrm{infinity})\,)\, $}

% \mapleresult
\begin{dmath}\label{(1.1.3)}
\mathit{MGF} \coloneqq \left(y ,c \right)\hiderel{\mapsto }\mathit{evalf} \!
\left(y -\left(\moverset{\infty}{\munderset{j =1}{\textcolor{gray}{\sum}}}\!
\frac{a \! \left(j ,c \right)}{y^{2\cdot j -1}}\right)\right)
\end{dmath}
\end{verbatim}
    \mapleinput
{$ \displaystyle \texttt{>\,} \mathit{MGF} \coloneqq (y ,c)\,\rightarrow
\,\mathit{evalf} (y -\mathit{Sum} (\frac{a (j ,c)}{y^{2\cdot j -1}},j =1..\mathrm{infinity})\,)\, $}

% \mapleresult
\begin{dmath}\label{(1.1.3)}
\mathit{MGF} \coloneqq \left(y ,c \right)\hiderel{\mapsto }\mathit{evalf} \! \left(y -\left(\moverset{\infty}{\munderset{j =1}{\textcolor{gray}{\sum}}}\! \frac{a \! \left(j ,c \right)}{y^{2\cdot j -1}}\right)\right)
\end{dmath}


\section{Figures}

The ``\verb|figure|'' environment should be used for figures. One or
more images can be placed within a figure. If your figure contains
third-party material, you must clearly identify it as such, as shown
in the example below.
\begin{figure}[h]
  \centering
  \includegraphics[width=\linewidth]{sample-franklin}
  \caption{1907 Franklin Model D roadster. Photograph by Harris \&
    Ewing, Inc. [Public domain], via Wikimedia
    Commons. (\url{https://goo.gl/VLCRBB}).}
  \Description{A woman and a girl in white dresses sit in an open car.}
\end{figure}

Your figures should contain a caption which describes the figure to
the reader.

Figure captions are placed {\itshape below} the figure.

Every figure should also have a figure description unless it is purely
decorative. These descriptions convey what’s in the image to someone
who cannot see it. They are also used by search engine crawlers for
indexing images, and when images cannot be loaded.

A figure description must be unformatted plain text less than 2000
characters long (including spaces).  {\bfseries Figure descriptions
  should not repeat the figure caption – their purpose is to capture
  important information that is not already provided in the caption or
  the main text of the paper.} For figures that convey important and
complex new information, a short text description may not be
adequate. More complex alternative descriptions can be placed in an
appendix and referenced in a short figure description. For example,
provide a data table capturing the information in a bar chart, or a
structured list representing a graph.  For additional information
regarding how best to write figure descriptions and why doing this is
so important, please see
\url{https://www.acm.org/publications/taps/describing-figures/}.

\subsection{The ``Teaser Figure''}

A ``teaser figure'' is an image, or set of images in one figure, that
are placed after all author and affiliation information, and before
the body of the article, spanning the page. If you wish to have such a
figure in your article, place the command immediately before the
\verb|\maketitle| command:
\begin{verbatim}
  \begin{teaserfigure}
    \includegraphics[width=\textwidth]{sampleteaser}
    \caption{figure caption}
    \Description{figure description}
  \end{teaserfigure}
\end{verbatim}



\section{Citations and Bibliographies}

The use of \BibTeX\ for the preparation and formatting of one's
references is \strikethrough{strongly recommended} \textcolor{red}{allowed}.
\textcolor{red}{Use of biblatex instead (an apparently more modern program) has some advantages, including the ``backref'' option. We recommend the style=numeric as being sufficiently similar to the ACM Recommended style for our purposes.}

\textcolor{red}{To use biblatex with back references (showing which page the work was cited on), which we recommend if you use the standard cite-as-you-go writing style, include the commands}
\begin{verbatim}
\usepackage[backref=true,style=numeric]{biblatex}
\addbibresource{sample-base.bib}
\end{verbatim}
\textcolor{red}{in the preamble, and instead of the bibliographystyle and bibliography commands at the end, use}
\begin{verbatim}
\printbibliography[title={Cited works with backrefs}]
\end{verbatim}
\textcolor{red}{(with a title that you want---the default is just ``References'') where you want the bibliography to be.  Several styles are available; consult the documentation for biblatex.}

\textcolor{red}{Citing online sources, such as the OEIS~\cite{sloane2021} or Wikipedia~\cite{wiki:ButcherGroup}, is encouraged.  The access date should be included in the .bib entry.}

Authors' names should be complete
--- use full first names (``Donald E. Knuth'') not initials
(``D. E. Knuth'') --- and the salient identifying features of a
reference should be included: title, year, volume, number, pages,
article DOI, etc.

The bibliography is included in your source document with these two
commands, placed just before the \verb|\end{document}| command:
\begin{verbatim}
  \bibliographystyle{ACM-Reference-Format}
  \bibliography{bibfile}
\end{verbatim}
where ``\verb|bibfile|'' is the name, without the ``\verb|.bib|''
suffix, of the \BibTeX\ file.
\textcolor{red}{Please use ``plain'' instead of ACM-Reference-Format for Maple Transactions. There is some incompatibility with some \LaTeX\ installations otherwise.  We hope this is temporary.}
\begin{verbatim}
  \bibliographystyle{plain}
  \bibliography{bibfile}
\end{verbatim}


Citations and references are numbered by default. \strikethrough{A small number of
ACM publications have citations and references formatted in the
``author year'' style; } $\ldots$
\textcolor{red}{We prefer that you not use ``author year'' style, but you may if you prefer.}



\textcolor{red}{Note: Please consider the practice of confining your citations to a single, final section, perhaps entitled ``Notes and References.''  This section should precede the acknowledgements which should immediately precede the references section generated by the bibliography command.  This practice (admittedly unusual in academia) can significantly enhance readability of the main text, while not compromising academic integrity.  If you choose to use this method, there will be little need to use backrefs.}


\section{Supplementary Material}
\textcolor{red}{Maple Transactions will publish ``live'' documents such as Maple Workbooks directly; if any code accompanies the paper, a link on the journal website next to the link to the paper will be provided.  In the text of the paper, it is still useful to provide a link to a repository where versions of the code can be found.  For instance, you could say something like the following:
Code supporting the paper can be found in
\href{https://github.com/rcorless/FractalEigenvector}{https://github.com/rcorless/FractalEigenvector}.
}

\section{Acknowledgments}

Identification of funding sources and other support, and thanks to
individuals and groups that assisted in the research and the
preparation of the work should be included in an acknowledgment
section, which is placed just before the reference section in your
document.

This section has a special environment:
\begin{verbatim}
  \begin{acks}
  ...
  \end{acks}
\end{verbatim}
so that the information contained therein can be more easily collected
during the article metadata extraction phase, and to ensure
consistency in the spelling of the section heading.

Authors should not prepare this section as a numbered or unnumbered {\verb|\section|}; please use the ``{\verb|acks|}'' environment.

\section{Appendices}

If your work needs an appendix, add it before the
``\verb|\end{document}|'' command at the conclusion of your source
document.

Start the appendix with the ``\verb|appendix|'' command:
\begin{verbatim}
  \appendix
\end{verbatim}
and note that in the appendix, sections are lettered, not
numbered. This document has two appendices, demonstrating the section
and subsection identification method.

\section{SIGCHI Extended Abstracts}

The ``\verb|sigchi-a|'' template style (available only in \LaTeX\ and
not in Word) produces a landscape-orientation formatted article, with
a wide left margin. Three environments are available for use with the
``\verb|sigchi-a|'' template style, and produce formatted output in
the margin:
\begin{itemize}
\item {\verb|sidebar|}:  Place formatted text in the margin.
\item {\verb|marginfigure|}: Place a figure in the margin.
\item {\verb|margintable|}: Place a table in the margin.
\end{itemize}

%%
%% The acknowledgments section is defined using the "acks" environment
%% (and NOT an unnumbered section). This ensures the proper
%% identification of the section in the article metadata, and the
%% consistent spelling of the heading.
\begin{acks}
\textcolor{red}{We thank the many authors of the \texttt{acmart} class for \LaTeX. That article class is very powerful, and its development shows considerable thought.  We hope that our derived class meets with their approval. We thank Eunice Y.~S.~Chan for her work; she was the one who did the alterations of \texttt{acmart} to create \texttt{MapleTrans}.  We thank Vic Martello and Kat McNichol for the Maple Transactions logo. We also thank Erik Postma for developing the Maple language set for the \texttt{listing} environment}.
\end{acks}

%%
%% The next two lines define the bibliography style to be used, and
%% the bibliography file.
\bibliographystyle{plain}
\bibliography{sample-base}
%\printbibliography

%%
%% If your work has an appendix, this is the place to put it.
\appendix

\section{Research Methods}

\subsection{Part One}

Lorem ipsum dolor sit amet, consectetur adipiscing elit. Morbi
malesuada, quam in pulvinar varius, metus nunc fermentum urna, id
sollicitudin purus odio sit amet enim. Aliquam ullamcorper eu ipsum
vel mollis. Curabitur quis dictum nisl. Phasellus vel semper risus, et
lacinia dolor. Integer ultricies commodo sem nec semper.

\subsection{Part Two}

Etiam commodo feugiat nisl pulvinar pellentesque. Etiam auctor sodales
ligula, non varius nibh pulvinar semper. Suspendisse nec lectus non
ipsum convallis congue hendrerit vitae sapien. Donec at laoreet
eros. Vivamus non purus placerat, scelerisque diam eu, cursus
ante. Etiam aliquam tortor auctor efficitur mattis.

\section{Online Resources}

Nam id fermentum dui. Suspendisse sagittis tortor a nulla mollis, in
pulvinar ex pretium. Sed interdum orci quis metus euismod, et sagittis
enim maximus. Vestibulum gravida massa ut felis suscipit
congue. Quisque mattis elit a risus ultrices commodo venenatis eget
dui. Etiam sagittis eleifend elementum.

% Nam interdum magna at lectus dignissim, ac dignissim lorem
% rhoncus. Maecenas eu arcu ac neque placerat aliquam. Nunc pulvinar
% massa et mattis lacinia.

\end{document}
\endinput
%%
%% End of file `sample-acmsmall.tex'.
